\section{Uniform recovery and spectral rigidity}
\label{sec:recovery}

The central case concerns a property whose natural bearer changes.  In a
Euclidean representation one proves decay of positive correlation forms.  In
a Hamiltonian representation the desired statement is absence of spectral
weight at low positive energy.  The bridge uses the OS reconstruction
paradigm \cite{OsterwalderSchrader1973,OsterwalderSchrader1975}, but the
abstract argument needs only positive semigroups and a dense recovered core.

\begin{definition}[Indexed semigroup recovery system]
\label{def:recovery}
Let $(I,\preceq)$ be directed.  For each $i\in I$, let $H_i\geq0$ be
self-adjoint on $\Hilb_i$.  Let $H\geq0$ be self-adjoint on $\Hilb$, with
distinguished unit vacuum $\vac$ and $\ker H=\mathbb C\vac$.  An indexed
semigroup recovery system consists of a linear subspace
$\mathcal D\subset\vac^\perp$ and eventually defined vectors
$J_i x\in\Hilb_i$ for $x\in\mathcal D$ such that, for every $t\geq0$,
\begin{equation}
\label{eq:semigroup-recovery}
 \langle J_i x,\e^{-tH_i}J_i x\rangle
 \longrightarrow
 \langle x,\e^{-tH}x\rangle.
\end{equation}
The system is \emph{complete} when $\mathcal D$ is dense in $\vac^\perp$.
It has \emph{uniform rigidity at rate $m$} when one $m>0$ exists such that,
for every $x\in\mathcal D$, a finite $M_x$ and index $i_x$ satisfy
\begin{equation}
\label{eq:uniform-rigidity}
 \langle J_i x,\e^{-tH_i}J_i x\rangle
 \leq M_x\e^{-mt}
 \quad (i\succeq i_x,\ t\geq0).
\end{equation}
\end{definition}

The common rate $m$ is uniform in the regulator index and is expressed in the
target physical time variable.  The prefactor may depend on the observable
vector, because spectral exclusion is established vector by vector before
density is used.  Equation~\eqref{eq:semigroup-recovery} may be obtained from
an explicit error modulus $\rho_i(x,t)$ that vanishes for every fixed $x,t$.

\begin{theorem}[Uniform recovery-to-spectrum]
\label{thm:recovery-spectrum}
Every complete indexed semigroup recovery system with uniform rigidity at
rate $m>0$ satisfies
\[
  \spec(H)\cap(0,m)=\varnothing.
\]
Equivalently, the target Hamiltonian has gap at least $m$ above its unique
vacuum.
\end{theorem}

\paragraph{Proof mechanism.}
For $x\in\mathcal D$, positivity and the spectral theorem give a positive
measure $\mu_x$ with
\[
 \langle x,\e^{-tH}x\rangle
   =\int_{[0,\infty)}\e^{-t\lambda}\,d\mu_x(\lambda).
\]
Taking the index limit in Equation~\eqref{eq:uniform-rigidity} gives an
$\e^{-mt}$ upper bound.  Any positive $\mu_x$-mass below $m$ would decay more
slowly and contradict that bound at large $t$.  Hence the spectral projection
$E_H((0,m))$ annihilates every $x\in\mathcal D$.  Density of $\mathcal D$ in
$\vac^\perp$ and boundedness of the projection imply
$E_H((0,m))=0$.  Appendix~\ref{app:spectral} gives the full proof and isolates
where positivity, centering, and density enter.

The theorem does not merely add representation errors.  It changes the type
of the property: a family of inequalities for recovered Euclidean quadratic
forms becomes a statement about the support of a target projection-valued
measure.  Positivity makes that transition possible; completeness makes it
global rather than sectorial; and uniformity prevents low-energy states from
escaping along the limiting index.

\begin{corollary}[OS-positive correlation-to-gap interface]
\label{cor:os}
Suppose a limiting Euclidean functional satisfies the hypotheses required for
OS reconstruction and yields $(\Hilb,\vac,H)$.  Suppose further that centered
vectors created by a gauge-invariant observable core are dense in
$\vac^\perp$, their regulated two-point forms converge to the reconstructed
semigroup forms, and Equation~\eqref{eq:uniform-rigidity} holds in physical
Euclidean time.  Then the reconstructed Hamiltonian has gap at least $m$.
\end{corollary}

This corollary is an interface theorem.  It does not prove OS positivity,
observable density, convergence of renormalized correlations, or the uniform
decay estimate for any particular interacting four-dimensional theory.

\subsection{A successful transport chain}

The elementary sequence
\[
 H_n=\operatorname{diag}(0,1,2,\ldots,n)
\]
embeds into $\mathbb C\vac\oplus\ell^2(\mathbb N)$ with target
$H\,e_k=k e_k$.  Every finite-support centered vector is recovered exactly
for all sufficiently large $n$; every regulated quadratic form decays at
least as $\e^{-t}$; finite-support vectors are dense; and the target gap is
one.  This example is simple, but it shows the complete logical chain:
carrier recovery, dynamical fidelity, a uniform rate, dense-core recovery,
and spectral interpretation.  The exact supplement checks its finite
truncations.  The counterexamples in the next section remove one certificate
at a time.
